
\chapter{Algunos scripts para el análisis de las trayectorias}


\begin{lstlisting}[language=Python, caption={Scripts de gromacs para remover condiciones periódicas, área por lípido, RMSD, perfil de densidad y radio de giro.}]
#!/bin/bash
mpi=/usr/bin/mpiexec
source /usr/local/gromacs/bin/GMXRC

# Set the number of MPI processes
export nprocs=6

#####Remover PBC del xtc
file=dopc_final

echo -e "ri 1-54\n ri 11-38 & a BB\n ri 65-110 & a BB | ri 117-120 & a BB\n 17 | 18\nq" | gmx_mpi make_ndx -f ${file}.tpr -o ${file}_rmsd.ndx
##Centrar todo en la caja
printf "System\nq"  |   gmx_mpi trjconv -f ${file}.xtc -s ${file}.tpr -n ${file}_rmsd.ndx -o ${file}_whole.xtc -pbc whole -tu ns

##Centrar la proteína usando uno de los péptidos
printf "r_1-54\nSystem\nq"  |   gmx_mpi trjconv -f ${file}_whole.xtc -s ${file}.tpr -n ${file}_rmsd.ndx -o ${file}_whole_centerResA.xtc -center -tu ns

###Centrar todo nuevamente (creo)
printf "System\nq"  |   gmx_mpi trjconv -f ${file}_whole_centerResA.xtc -s ${file}.tpr -n ${file}_rmsd.ndx -o ${file}_whole_centerResA_mol_noPBC.xtc -pbc mol -tu ns -ur compact

rm ${file}_whole.xtc ${file}_whole_centerResA.xtc
mkdir -p "./results/rmsd/" && mv ${file}_rmsd* "./results/rmsd/" 

####RMSD #Protein BB
printf "19\n19\nq" |   gmx_mpi rms -s ${file}.tpr -f ${file}_whole_centerResA_mol_noPBC.xtc -o ${file}_rmsdProtein.xvg -n  ./results/rmsd/${file}_rmsd.ndx -tu ns

#Alfa BB
printf "17\n17\nq" |   gmx_mpi rms -s ${file}.tpr -f ${file}_whole_centerResA_mol_noPBC.xtc -o ${file}_rmsdProA.xvg -n  ./results/rmsd/${file}_rmsd.ndx -tu ns

#Beta BB
printf "18\n18\nq" |   gmx_mpi rms -s ${file}.tpr -f ${file}_whole_centerResA_mol_noPBC.xtc -o ${file}_rmsdProB.xvg -n  ./results/rmsd/${file}_rmsd.ndx -tu ns

mv ${file}_rmsd* './results/rmsd/'

###edr para APL
  gmx_mpi eneconv -f dopc.edr dopc_ext1.edr dopc_ext2.edr dopc_ext3.edr dopc_ext4.edr dopc_ext5.edr -o ${file}.edr

printf "Box-X\nBox-Y\nSystem\n" |   gmx_mpi energy -f ${file}.edr -o ${file}_ApL.xvg

mkdir -p "./results/ApL_Density/" &&
mv ${file}_ApL.xvg "./results/ApL_Density/"

####Density profile
echo -e "a PO4 a NC3\n a GL1 a GL2\n13 & ! 16 & ! 17 & ! a NC3 \nq" |   gmx_mpi make_ndx -f ${file}.tpr -o ${file}_density_groups.ndx

## -center usar "0" (System) como refencia
printf "13\n16\n" |   gmx_mpi density -s ${file}.tpr -f ${file}_whole_centerResA_mol_noPBC.xtc -n ${file}_density_groups.ndx -o ${file}_den_headgroups.xvg -d Z -center -sl 100
printf "13\n17\n" |   gmx_mpi density -s ${file}.tpr -f ${file}_whole_centerResA_mol_noPBC.xtc -n ${file}_density_groups.ndx -o ${file}_den_glycerol_ester.xvg -d Z -center -sl 100
printf "13\n18\n" |   gmx_mpi density -s ${file}.tpr -f ${file}_whole_centerResA_mol_noPBC.xtc -n ${file}_density_groups.ndx -o ${file}_den_acylChain.xvg -d Z -center -sl 100
printf "13\n14\n" |   gmx_mpi density -s ${file}.tpr -f ${file}_whole_centerResA_mol_noPBC.xtc -n ${file}_density_groups.ndx -o ${file}_den_water.xvg -d Z -center -sl 100

###Crear carpeta y mover todos los files recién creados
mkdir -p "./results/ApL_Density"
mv ${file}_den* "./results/ApL_Density"

###Rgiro
printf "Protein\nSystem\n" |   gmx_mpi gyrate -s ${file}.tpr -f ${file}_whole_centerResA_mol_noPBC.xtc -o ${file}_radius_g.xvg

mkdir -p "./results/radius_g"
mv ${file}_radius* "./results/radius_g"
\end{lstlisting}



\begin{lstlisting}[language=Python, caption=Script para calcular mapa de contactos péptido-péptido.]
#!/usr/bin/env python3
import MDAnalysis as mda
import numpy as np
import MDAnalysis.analysis.distances
from tqdm.auto import tqdm
from pathlib import Path
from mpi4py import MPI

# Initialize MPI
comm = MPI.COMM_WORLD
rank = comm.Get_rank()
size = comm.Get_size()

XTC = 'dppc_noPBC.xtc'  # xtc periodic conditions removed
TPR = 'dppc.pdb'
u = mda.Universe(TPR, XTC)

def contact_matrix(protein_1, protein_2, start_time, end_time):
    contacts = np.zeros((len(np.unique(protein_1.resids)), len(np.unique(protein_2.resids))))

    # Determine frame indices corresponding to start and end times
    total_frames = len(u.trajectory)
    start_frame = int(start_time / 4.5 * total_frames)
    end_frame = int(end_time / 4.5 * total_frames)

    # Get sub_1 and sub_2
    sub_1 = min(np.unique(protein_1.resids))
    sub_2 = min(np.unique(protein_2.resids))

    # Iterate over trajectory frames within the specified range
    for ts in tqdm(u.trajectory[start_frame:end_frame]):
        if ts.frame % size == rank:
            for i in np.unique(protein_1.resids):
                a = protein_1.select_atoms('resid '+str(i))
                for k in np.unique(protein_2.resids):
                    b = protein_2.select_atoms('resid '+str(k))
                    if MDAnalysis.analysis.distances.distance_array(a.center_of_mass(), b.center_of_mass(), box=[u.dimensions[0], u.dimensions[1], u.dimensions[2], u.dimensions[3], u.dimensions[4], u.dimensions[5]]) < 7.0:
                        contacts[i-sub_1, k-sub_2] += 1

    return contacts

# Select sequence/proteins
protein_1 = u.select_atoms('index 0-135')
protein_2 = u.select_atoms('index 136-350')

# Define time intervals
time_intervals = [
    (0, 0.5),
    (0.5, 1.5),
    (1.5, 2.5),
    (2.5, 3.5),
    (3.5, 4.5)
]

# Run function for each time interval
for idx, (start_time, end_time) in enumerate(time_intervals):
    if rank == 0:
        print(f"Processing time interval {idx+1}: start_time={start_time}, end_time={end_time}")
    contact_matrix_result = contact_matrix(protein_1, protein_2, start_time, end_time)

    # Gather results from all processes
    all_contact_matrices = comm.gather(contact_matrix_result, root=0)

    # Save results on root process
    if rank == 0:
        results_directory = Path("results/contacts")
        results_directory.resolve().mkdir(exist_ok=True, parents=True)
        np.save(results_directory / f'dppc_{start_time}_{end_time}us.npy', np.sum(all_contact_matrices, axis=0))

\end{lstlisting}


\begin{lstlisting}[language=Python, caption=Script para graficar mapa de contactos péptido-péptido.]
#!/usr/bin/env python3
import numpy as np
import matplotlib.pyplot as plt
import seaborn as sns
import pandas as pd
from matplotlib.ticker import MultipleLocator  # Import MultipleLocator
from matplotlib.ticker import AutoMinorLocator
import matplotlib.ticker as ticker
import MDAnalysis as mda


##Load here the trajectory and pdb files
XTC = 'dppc_noPBC.xtc'  # xtc periodic conditions removed
TPR = 'dppc.pdb'
u = mda.Universe(TPR, XTC)

###Function to plot
def set_plot_style():
    plt.rcParams["figure.figsize"] = (6, 4)
    plt.rcParams['font.family'] = 'serif'
    plt.rcParams['font.sans-serif'] = ['Arial']
    plt.rcParams.update({'font.size': 24})
    plt.rc('font', weight='bold')
    plt.rcParams['figure.autolayout'] = True
    plt.rcParams['image.interpolation'] = 'nearest'
    plt.rcParams['axes.unicode_minus'] = False
    plt.rcParams['figure.dpi'] = 600
    plt.rcParams["ytick.labelsize"] = 15
    plt.rcParams["xtick.labelsize"] = 15
    plt.rcParams["axes.linewidth"] = 2.0     # grosor del frame box
    plt.rcParams['xtick.major.width'] = 2.0  # grosor de indicar número en x
    plt.rcParams['ytick.major.width'] = 2.0  # grosor de indicar número en y
    plt.rcParams['ytick.minor.visible'] = True
    plt.rcParams['ytick.minor.width'] = 2.0
    plt.rcParams['xtick.minor.width'] = 2.0
    plt.rcParams['xtick.minor.visible'] = True
    plt.rc('xtick.minor', size=3, pad=1)
    plt.rc('ytick.minor', size=3, pad=1)
    plt.rcParams["axes.labelsize"] = 18
    plt.rcParams["axes.labelcolor"] = 'black'
    plt.rcParams["axes.labelweight"] = 'bold'
    plt.rcParams["axes.titlesize"] = 10
    plt.rcParams["axes.titleweight"] = 'bold'
    plt.rcParams['savefig.dpi'] = 600
    plt.rcParams['savefig.bbox'] = 'tight'
    plt.rcParams['savefig.pad_inches'] = 0.05
    plt.rcParams["legend.fontsize"] = 18
#     matplotlib.rcParams['axes.linewidth'] = 3


# Load labels
y_label_beta = np.load('results/contacts/y_label_beta.npy', allow_pickle=True)
x_label_alfa = np.load('results/contacts/x_label_alfa.npy', allow_pickle=True)

# Loop over each file. This for rename xlabel in all matrix with residues Leter+Number
for file_name in [file+'_0_0.5us.npy', file+'_0.5_1.5us.npy', file+'_1.5_2.5us.npy', 
                  file+'_2.5_3.5us.npy', file+'_3.5_4.5us.npy']:
    # Load contacts matrix
    contacts_mat = np.load(f'results/contacts/{file_name}') / len(u.trajectory)
    
    # Normalize contacts matrix
    contacts_norma = (contacts_mat - contacts_mat.min()) / (contacts_mat.max() - contacts_mat.min())
    
    # Create DataFrame
    df = pd.DataFrame(contacts_norma, columns=[y_label_beta])
    df.index = [x_label_alfa]
    
    # Save DataFrame to CSV
    csv_file_name = file_name.replace('.npy', '.csv')
    df.to_csv(f'results/contacts/contactMap_{csv_file_name}')


# Load labels
y_label_beta = np.load('results/contacts/y_label_beta.npy', allow_pickle=True)
x_label_alfa = np.load('results/contacts/x_label_alfa.npy', allow_pickle=True)

###Lineas para mostrar contactos        
contact_points = [
    ("G972", 16.5, 17.5),
    ("G976", 20.5, 21.5),
    ("L979", 24.5, 24.5),
    ("R995", 39.5, 40.5)]

contact_points2 = [("M987", 35.5, 32.5)]
contact_points3 = [("R995", 39.5, 40.5)]

# Define file names
files = [
    'contactMap_'+file+'_0_0.5us.csv',
    'contactMap_'+file+'_0.5_1.5us.csv',
    'contactMap_'+file+'_1.5_2.5us.csv',
    'contactMap_'+file+'_2.5_3.5us.csv',
    'contactMap_'+file+'_3.5_4.5us.csv'
]

set_plot_style()

# Calculate number of rows and columns based on the number of files
num_files = len(files)
num_cols = 2
num_rows = (num_files + num_cols - 1) // num_cols  # Ceiling division

# Create subplots
fig, axes = plt.subplots(figsize=(25,8 * num_rows), nrows=num_rows, ncols=num_cols, sharex=False, sharey=False, 
                         gridspec_kw={'wspace': 0.15, 'hspace': 0.35})

# Loop over each file and plot the heatmap
for i, file_name in enumerate(files):
    df = pd.read_csv(f'results/contacts/{file_name}', header=0, index_col=0)
    row = i // num_cols
    col = i % num_cols
    sns.heatmap(df, cmap="Greys", ax=axes[row, col], cbar_kws={'shrink': 0.99, "pad": 0.15, "label": 'Frecuencia normalizada'})

    # Customize axes
    axes[row, col].set_xlabel(r"$\beta$3", fontsize=30, color='blue')
    axes[row, col].set_ylabel(r"$\alpha$IIb", fontsize=30, color='red')
    axes[row, col].xaxis.set_minor_locator(plt.NullLocator())
    axes[row, col].yaxis.set_minor_locator(plt.NullLocator())
    
    axes[row, col].spines['top'].set_visible(True)
    axes[row, col].spines['right'].set_visible(True)
    axes[row, col].spines['bottom'].set_visible(True)
    axes[row, col].spines['left'].set_visible(True)
    
    # Set x-ticks every 5 labels
    n = 5
    xticks = range(0, len(df.columns), n)
    xticklabels = df.columns[::n]
    axes[row, col].set_xticks(xticks)
    axes[row, col].set_xticklabels(xticklabels, fontsize=20)
    
    # Set y-ticks every 5 labels
    yticks = range(0, len(df.index), n)
    yticklabels = df.index[::n]
    axes[row, col].set_yticks(yticks)
    axes[row, col].set_yticklabels(yticklabels, fontsize=20)   

    # Set the x and y limits
    axes[row, col].set_xlim(10, 45)  # Beta
    axes[row, col].set_ylim(10, 45)  # Alfa
    
    axes[row, col].invert_yaxis()
    
    # Add text
    if i == 0:
        axes[row, col].text(0.5, 0.95, r'0.5 $\mu$s', horizontalalignment='center', verticalalignment='top', 
                            transform=axes[row, col].transAxes, fontsize=30)
    elif i == 1:
        axes[row, col].text(0.5, 0.95, r'1.5 $\mu$s', horizontalalignment='center', verticalalignment='top', 
                            transform=axes[row, col].transAxes, fontsize=30)
    elif i == 2:
        axes[row, col].text(0.5, 0.95, r'2.5 $\mu$s', horizontalalignment='center', verticalalignment='top', 
                            transform=axes[row, col].transAxes, fontsize=30)
    elif i == 3:
        axes[row, col].text(0.5, 0.95, r'3.5 $\mu$s', horizontalalignment='center', verticalalignment='top', 
                            transform=axes[row, col].transAxes, fontsize=30)
    elif i == 4:
        axes[row, col].text(0.5, 0.95, r'4.5 $\mu$s', horizontalalignment='center', verticalalignment='top', 
                            transform=axes[row, col].transAxes, fontsize=30)
     
    # Add contact lines
    for contact, x1, y1 in contact_points:
        axes[row, col].plot([x1, 79], [y1, y1], color='m', ls=':', lw=3.0)  # Horizontal line
        axes[row, col].plot([x1, x1], [y1, 79], color='m', ls=':', lw=3.0)  # Vertical line
    
    for contact, x1, y1 in contact_points2:
        axes[row, col].plot([x1, 79], [y1, y1], color='#FE5100', ls=':', lw=3.0)  # Horizontal line
        axes[row, col].plot([x1, x1], [y1, 79], color='#FE5100', ls=':', lw=3.0)  # Vertical line  
    
    for contact, x1, y1 in contact_points3:
        axes[row, col].plot([x1, 79], [y1, y1], color='blue', ls=':', lw=3.0)  # Horizontal line
        axes[row, col].plot([x1, x1], [y1, 79], color='blue', ls=':', lw=3.0)  # Vertical line

    # Add additional text for contacts
    axes[row, col].text(45.5, 18, 'G972', fontsize=20, color='m')
    axes[row, col].text(15.9, 50.9, 'V700', fontsize=20, color='m', rotation=90)
    
    axes[row, col].text(45.5, 22.5, 'G976', fontsize=20, color='m')
    axes[row, col].text(19.24, 50.50, 'I704', fontsize=20, color='m', rotation=90)
    
    axes[row, col].text(45.5, 24.5, 'L979', fontsize=20, color='m')
    axes[row, col].text(22.5, 51.3, 'G708', fontsize=20, color='m', rotation=90)
    
    axes[row, col].text(45.5, 33.3, 'M987', fontsize=20, color='#FE5100')
    axes[row, col].text(34.25, 50.498, 'I719', fontsize=20, color='#FE5100', rotation=90)
    
    axes[row, col].text(45.62, 40.7, 'R995', fontsize=20, color='b')
    axes[row, col].text(38, 51.3, 'D723', fontsize=20, color='b', rotation=90)
    

# Turn off minor ticks in the colorbar
for ax in fig.get_axes():
    if isinstance(ax, plt.Axes):
        ax.xaxis.set_minor_locator(plt.NullLocator())
        ax.yaxis.set_minor_locator(plt.NullLocator())

# Remove empty subplot if present
if num_files % num_cols != 0:
    fig.delaxes(axes[num_files // num_cols, num_cols - 1])

# Colocar segundo Yaxis a la derecha
for ax_row in axes:
    for ax in ax_row:
        ax.tick_params(axis='y', right=True, labelright=True, which='both', rotation=0)
        
        
fig.text(0.7, 0.2, "DPPC", ha='center', va='center',fontsize=80)

plt.savefig('results/contacts/contactMap_'+file+'.pdf')
plt.savefig('results/contacts/contactMap_'+file+'.svg')

plt.show()
\end{lstlisting}



\begin{lstlisting}[language=Python,caption={Script para calcular ángulo de inclinación}]
#!/bin/bash
source /usr/local/gromacs/bin/GMXRC
# Define the time intervals
time_intervals=(
    "0 500000"
    "500000 1500000"
    "1500000 2500000"
    "2500000 3500000"
    "3500000 4500000"
)

file=dppc
# Create an index for alfa and beta
echo -e "ri 12-39 & a BB\ndel 0-15\nq" | gmx make_ndx -f ${file}.tpr -o ${file}_index_tilt_alfa.ndx
echo -e "ri 65-92 & a BB\ndel 0-15\nq" | gmx make_ndx -f ${file}.tpr -o ${file}_index_tilt_beta.ndx

# Loop over each time interval and run gmx principal
for interval in "${time_intervals[@]}"; do
    start=$(echo "$interval" | cut -d' ' -f1)
    end=$(echo "$interval" | cut -d' ' -f2)

    echo "Processing time interval: ${start}-${end}"
    
    printf "${start}\n${end}\nq" | gmx principal -f ${file}_noPBC.xtc -s ${file}.tpr -n ${file}_index_tilt_alfa.ndx -a1 paxis1_alfa_${start}_${end}.xvg -a2 paxis2_alfa_${start}_${end}.xvg -a3 paxis3_alfa_${start}_${end}.xvg -om omi_alfa_${start}_${end}.xvg -b ${start} -e ${end}
    printf "${start}\n${end}\nq" | gmx principal -f ${file}_noPBC.xtc -s ${file}.tpr -n ${file}_index_tilt_beta.ndx -a1 paxis1_beta_${start}_${end}.xvg -a2 paxis2_beta_${start}_${end}.xvg -a3 paxis3_beta_${start}_${end}.xvg -om omi_beta_${start}_${end}.xvg -b ${start} -e ${end}

    echo "Completed processing for interval: ${start}-${end}"
done

# Move all created files to the "./results/tilt_angle" directory
mkdir -p "./results/tilt_angle" && mv ${file}_index_tilt* paxis* omi* "./results/tilt_angle"
\end{lstlisting}

%\newpage

\begin{lstlisting}[language=Python, caption=Script para graficar ángulos de inclinación.]
#!/usr/bin/env python3
import numpy as np
import matplotlib.pyplot as plt
import seaborn as sns
import pandas as pd
from matplotlib.ticker import MultipleLocator  # Import MultipleLocator
from matplotlib.ticker import AutoMinorLocator
import matplotlib.ticker as ticker
import warnings

# Define the time intervals 
time_intervals = [
    (0, 500000),
    (500000, 1500000),
    (1500000, 2500000),
    (2500000, 3500000),
    (3500000, 4500000)]

set_plot_style() ###should be defined before
# Calculate number of rows and columns based on the number of files
num_files = 5
num_cols = 2
num_rows = (num_files + num_cols - 1) // num_cols  # Ceiling division

fig, axes = plt.subplots(figsize=(14,5 * num_rows), nrows=num_rows, ncols=num_cols, 
                         gridspec_kw={'wspace': 0.25, 'hspace': 0.3})

for idx, (start, end) in enumerate(time_intervals):
    file_path = f'results/tilt_angle/'

    # Load data with warning suppression
    with warnings.catch_warnings():
        warnings.simplefilter("ignore")  # Suppress conversion warnings
        data1 = np.genfromtxt([i for i in open(f'{file_path}paxis1_alfa_{start}_{end}.xvg').read().splitlines() 
                               if not i.startswith(('#','@'))]) 
        data2 = np.genfromtxt([i for i in open(f'{file_path}paxis1_beta_{start}_{end}.xvg').read().splitlines() 
                               if not i.startswith(('#','@'))])

        # Check if there are any errors in loading the data
        if np.isnan(data1).any() or np.isnan(data2).any():
            print(f"There were errors loading data for interval {start}-{end}.")
        else:

            # Calculate tilt angles
            df_alfa = pd.DataFrame(data1, columns=['Time_ps', 'x_component', 'y_component', 'z_component'])
            df_alfa['magnitude'] = np.sqrt(df_alfa['x_component']**2 + df_alfa['y_component']**2 + df_alfa['z_component']**2)
            df_alfa['tilt_ang'] = (np.arccos(df_alfa['z_component']) / df_alfa['magnitude']) * 180 / np.pi

            df_beta = pd.DataFrame(data2, columns=['Time_ps', 'x_component', 'y_component', 'z_component'])
            df_beta['magnitude'] = np.sqrt(df_beta['x_component']**2 + df_beta['y_component']**2 + df_beta['z_component']**2)
            df_beta['tilt_ang'] = (np.arccos(df_beta['z_component']) / df_beta['magnitude']) * 180 / np.pi

            # Calculate crossing angle
            df_alfa['cross_ang'] = np.arccos((df_alfa['x_component']*df_beta['x_component'] + df_alfa['y_component']*df_beta['y_component'] + df_alfa['z_component']*df_beta['z_component']) / (df_alfa['magnitude']*df_beta['magnitude'])) * 180 / np.pi

            # Save to CSV
            df_alfa.to_csv(f'{file_path}dppc_{start}_{end}.csv')
            
            row = idx // 2
            col = idx % 2
            ax = axes[row, col]
            sns.distplot(df_alfa['tilt_ang'], hist=True, kde=True, bins=20, color='red', label=r'$\alpha$IIb' f"\nMean: {df_alfa['tilt_ang'].mean():.2f} ± {df_alfa['tilt_ang'].std():.2f}\textdegree ", ax=ax)
            sns.distplot(df_beta['tilt_ang'], hist=True, kde=True, bins=60, color='blue', label=r'$\beta$3' f"\nMean: {df_beta['tilt_ang'].mean():.2f} ± {df_beta['tilt_ang'].std():.2f}\textdegree ", ax=ax)
            sns.distplot(df_alfa['cross_ang'], hist=True, kde=True, bins=80, color='darkgreen', label=r'$\alpha$IIb$\beta$3' f"\nMean: {df_alfa['cross_ang'].mean():.2f} ± {df_alfa['cross_ang'].std():.2f}\textdegree ", ax=ax)

            ax.set_xlabel('Angle (\textdegree )')  # Plot formatting
            ax.set_ylabel('Density')
            ax.set_xlim(0, 60)
            ax.set_ylim(0, 0.180)
            ax.legend(loc=2, frameon=False, fontsize=12)
            ax.set_title([f' 0.5 $\mu$s', f' 1.5 $\mu$s', f' 2.5 $\mu$s', f' 3.5 $\mu$s', f' 4.5 $\mu$s'][idx], fontsize=18)

# Remove empty subplot if present
if len(time_intervals) % 2 != 0:
    fig.delaxes(axes[len(time_intervals) // 2, 1])
    
# Turn off minor ticks in the colorbar
for ax in fig.get_axes():
    if isinstance(ax, plt.Axes):
        ax.xaxis.set_minor_locator(plt.NullLocator())
        ax.yaxis.set_minor_locator(plt.NullLocator())   
    
fig.text(0.7, 0.2, "DPPC", ha='center', va='center',fontsize=40)

plt.savefig('results/tilt_angle/dppcAngles.pdf')
plt.savefig('results/tilt_angle/dppcAngles.svg')
plt.show()
\end{lstlisting}


\begin{lstlisting}[language=Python, caption={Script para calcular las frecuencias de interacción lípido-péptido.}]
#!/usr/bin/env python3
import MDAnalysis as mda
import numpy as np
import pandas as pd
import matplotlib.pyplot as plt
from MDAnalysis.analysis.leaflet import LeafletFinder
from tqdm.auto import tqdm

# Load an MDAnalysis Universe
XTC = 'dppc_noPBC.xtc' # xtc periodic conditions removed
TPR = 'dppc.pdb'
u = mda.Universe(TPR, XTC, topology_format='PDB')

####Para asignar seg############
# list(u.select_atoms('protein')) #del 1 al 136 es alfa y del 137 al 351 es beta
alpha = u.add_Segment(segid='seg_0_PROA_P')
alpha_atoms = u.select_atoms('protein and bynum 1:136')
alpha_atoms.residues.segments = alpha
alpha_atoms.atoms

beta = u.add_Segment(segid='seg_1_PROB_P')
beta_atoms = u.select_atoms('protein and bynum 137:351')
beta_atoms.residues.segments = beta
beta_atoms.atoms
#####FIN############


###Seleccionamos solo el TMD
prot_A = u.select_atoms('protein and segid seg_0_PROA_P and (resid 9-40)')
# list(prot_A)

# Para normalizar respecto al n total de LIPIDS
num_LIPIDS = len(u.select_atoms("resname DPPC"))

# Define selections for protein and LIPID
protein = prot_A
lipid = u.select_atoms("resname DPPC and name PO4")

# Create an array to store contact frequencies for each residue
contact_frequencies = np.zeros(len(protein.resids))

# Loop through all frames
for ts in tqdm(u.trajectory[::]):  # para hacer prueba cada mil frame [::1000]
    # Calculate distances between protein and lipid atoms
    distances = np.linalg.norm(protein.positions[:, np.newaxis] - lipid.positions, axis=2)

    # Check if any distance is below the cutoff
    contacts_matrix = distances < 6.0

    # Sum the contacts for each residue
    contact_residues = np.unique(protein.resids[np.where(contacts_matrix.any(axis=1))[0]])
    contact_frequencies[contact_residues - 1] += 1  # Subtract 1 to match array indices

# Normalize the contact frequencies by the number of lipids
contact_frequencies_normalized = contact_frequencies / num_LIPIDS

# Create a range for residue numbers starting from 1
residue_numbers = range(1, len(protein.resids) + 1)

# Create a DataFrame to hold the data
data = pd.DataFrame({'residue_number': residue_numbers, 'normalized_contact_frequency': contact_frequencies_normalized})

###Save file
from pathlib import Path
results_directory = Path("results/lipid_protein_contact")
results_directory.resolve().mkdir(exist_ok=True, parents=True)

csv_file_path = results_directory / "dppc_alpha_contact_normalized.csv"
# Save the DataFrame as a CSV file in the results directory
data.to_csv(csv_file_path, index=False)

plt.bar(residue_numbers, contact_frequencies_normalized, align='center')
plt.xlabel(r'$\alpha$IIb Residue Number')
plt.ylabel('Contact Frequency')
plt.title('Contact Frequencies for Protein Residues')
plt.xlim(0,54)
plt.show()
\end{lstlisting}



\begin{lstlisting}[language=Python, caption={Script para calcular la ocupancia de los lípidos. Definida como porcentaje de frames en la trayectoria donde los lípidos estuvieron a una distancia <\,6.0 \AA \; respecto a algún residuo de los péptidos.}]
import MDAnalysis as mda
import numpy as np
import pandas as pd
from tqdm.auto import tqdm

# Load an MDAnalysis Universe
XTC = 'dppc_noPBC.xtc'
TPR = 'dppc.tpr'
u = mda.Universe(TPR, XTC)

# Select Proteins
prot_A = u.select_atoms('protein and segid seg_0_PROA_P and (resid 9-40)')
prot_B = u.select_atoms('protein and segid seg_1_PROB_P and (resid 63-94)')

# Select Lipids
lipid = u.select_atoms("resname POPC and name PO4")

# Initialize dictionaries to store contact counts for each lipid
contact_counts_A = {lipid_index: 0 for lipid_index in range(len(lipid))}
contact_counts_B = {lipid_index: 0 for lipid_index in range(len(lipid))}

# Define a cutoff for contact (in Ångstroms)
cutoff_distance = 6.0

# Iterate through frames
total_frames = 0
for ts in tqdm(u.trajectory[::5]):
    total_frames += 1
    distances_A = np.linalg.norm(prot_A.positions[:, np.newaxis] - lipid.positions, axis=2)
    contacts_matrix_A = distances_A < cutoff_distance
    
    distances_B = np.linalg.norm(prot_B.positions[:, np.newaxis] - lipid.positions, axis=2)
    contacts_matrix_B = distances_B < cutoff_distance

    # Update contact count for each lipid for prot_A and prot_B
    for lipid_index in range(len(lipid)):
        if any(contacts_matrix_A[:, lipid_index]):
            contact_counts_A[lipid_index] += 1
        if any(contacts_matrix_B[:, lipid_index]):
            contact_counts_B[lipid_index] += 1

# Calculate occupancy as a percentage
alfa_occupancy = {lipid_index: (count / total_frames) * 100 for lipid_index, count in contact_counts_A.items()}
beta_occupancy = {lipid_index: (count / total_frames) * 100 for lipid_index, count in contact_counts_B.items()}

# Create DataFrame
occupancy_df = pd.DataFrame({
    'Lipid_Index': list(alfa_occupancy.keys()),
    'Alfa_Occupancy': list(alfa_occupancy.values()),
    'Beta_Occupancy': list(beta_occupancy.values())
})

# Print DataFrame to check
occupancy_df.to_csv('results/lipid_protein_contact/dppc_ocupancy.csv', index=False)
occupancy_df.head()


###Plot
occupancy_df = pd.read_csv('results/lipid_protein_contact/popc_occupancy.csv')

# Melt the DataFrame to have a long-form DataFrame which seaborn can use to plot
occupancy_melted_df = occupancy_df.melt(id_vars='Lipid_Index', value_vars=['Alfa_Occupancy', 'Beta_Occupancy'],
                                        var_name='Selection', value_name='Occupancy')

# Create the violin plot
plt.figure(figsize=(10, 6))
sns.violinplot(x='Selection', y='Occupancy', data=occupancy_melted_df)

plt.title('Lipid Occupancy Near Protein Alfa and Beta Selections')
plt.xlabel('Protein Selection')
plt.ylabel('Occupancy (%)')
plt.show()

\end{lstlisting}
